% plan_v2_vn.tex -- the plan for paper_v2_vn: the Vietnamese version of paper_v2. A standalone
% document; compile it with LuaLaTeX, like paper_v2_vn.tex (Section "Technical setup"). The
% folders paper/ and paper_v2/ are not edited for this work (Section "Rules").
\documentclass[11pt,a4paper]{article}

\usepackage{iftex}
\RequireLuaTeX   % Vietnamese letters need LuaLaTeX here, as in paper_v2_vn.tex
\usepackage{fontspec}   % Latin Modern, which has every Vietnamese letter
\usepackage[vietnamese,main=english]{babel}
\usepackage[margin=2.2cm]{geometry}
\usepackage{microtype}
\usepackage{icomma}   % a comma between digits in a formula is a decimal comma
\usepackage{array}
\usepackage{booktabs}
\usepackage{tabularx}
\usepackage{enumitem}
\usepackage[hidelinks]{hyperref}

\DeclareRobustCommand{\code}[1]{\texttt{\detokenize{#1}}}
\newcommand{\vi}[1]{\foreignlanguage{vietnamese}{#1}}
\newcolumntype{L}{>{\raggedright\arraybackslash}X}
\setlist{itemsep=2pt,topsep=3pt}
\setlength{\parskip}{4pt}
\setlength{\parindent}{0pt}
\setlength{\emergencystretch}{3em}

\title{Plan for \texttt{paper\_v2\_vn}\\[4pt]\large The Vietnamese version of \texttt{paper\_v2}}
\author{}
\date{10 November 2026}

\begin{document}
\maketitle

\section*{Status}

All nine parts are done (10 November): \code{paper_v2_vn/} is the whole paper in Vietnamese and
builds with LuaLaTeX with 0 warnings (Section~\ref{sec:l9} gives the final checks). The author
committed \code{make_numbers.py} as \code{5e32eb6} of the LMSPS repository the same day, and
\code{numbers.tex}, rebuilt from it, records that commit; no other value changed. One step is
left, the author's: reading the Vietnamese text, ideally with the supervisors, since the AI
statement says that the author reviewed the translation; corrections from that reading come in as
L10, which started on 20 November (N17). The author's
answers are decisions N3--N32 (Section~\ref{sec:decisions}); the parts are in
Section~\ref{sec:parts}.

\section{The target}

A Vietnamese version of \code{paper_v2} with the same content, numbers and layout: ACM's
two-column format (\code{acmart}, \code{sigconf}, \code{nonacm}) in its own Libertine fonts, and
every claim, qualifier, number, citation, table and figure of \code{paper_v2}, with nothing added
and nothing dropped. Vietnamese text usually takes somewhat more room than the English it
translates, so the main text will probably end half a page to a page and a half after
\code{paper_v2}'s 8 pages. Up to about one page more is fine (N3); L9 measures it.

\section{Rules}\label{sec:rules}

\begin{itemize}
  \item \textbf{\code{paper/} and \code{paper_v2/} are frozen.} They change only for a must: a
        factual error found while translating (a wrong number, claim, reference or
        cross-reference). I report such an error first; its fix then goes into both folders, as
        decision V5 of \code{plan_v2} did, and into \code{paper_v2_vn}. Wording that could
        merely be better stays as it is.
  \item \textbf{A faithful translation.} \code{paper_v2} is the source. The claims keep their
        qualifiers (``mostly'', seed noise, the warm-up limit, the small schema), and no claim
        is added or dropped. Where the English could be read two ways, I ask rather than guess.
  \item \textbf{Every number comes from \code{numbers.tex}} as \verb|\val{key}|. A script
        makes that file and nobody types it (N4). Design
        constants typed in the text, such as 80/10/10, follow the same number format.
  \item \textbf{What stays in English:} the names of models (HAN, GAT, LMSPS, metapath2vec),
        metrics (Hit@10, NDCG@10, BPR), paths (\code{PUP}, \code{PCCP}), the dataset
        (Amazon-Meta), files, commands and code, and the reference list.
  \item \textbf{The AI statement stays accurate.} The Vietnamese one adds that the Vietnamese
        text was translated from the English with the AI assistant's help and that the author
        reviewed the translation. So the author has to review it, ideally with the supervisors,
        before the version is used.
  \item \textbf{Each part ends with checks:} the translated files carry exactly the same
        \verb|\val|, \verb|\cite|, \verb|\ref|, \verb|\label|, \verb|\input| and \verb|\code|
        keys as \code{paper_v2}'s (checked by script); a LuaLaTeX build with 0 warnings; and a
        page count. It sends only
        the files it added or changed (L1 sends the new folder once).
\end{itemize}

\section{The folder}\label{sec:folder}

\code{paper_v2_vn/} mirrors \code{paper_v2/}, with the same file names so that each file can be
checked against its source: \code{paper_v2_vn.tex}, \code{numbers.tex}, \code{sections/},
\code{tables/}, \code{figures/}, this plan and the glossary (\code{glossary_vn.tex}, L2). The three \code{.bib} files stay at the project root,
shared by all three versions; the reference list stays in English. Since N31 the folder also holds
the supplementary document, \code{supplement_vn.tex}, built after the paper (\code{latexmkrc}
builds both).

\section{Technical setup (tested 10 November)}\label{sec:setup}

\begin{itemize}
  \item \textbf{pdfLaTeX cannot do it.} Under pdfLaTeX, \code{acmart} uses Libertine fonts that
        have no Vietnamese encoding (T5), so Vietnamese letters fail (``Font shape
        \code{T5/LinuxLibertineT-TLF/m/n} undefined'').
  \item \textbf{LuaLaTeX can.} Under LuaLaTeX, \code{acmart} loads Libertine's OpenType fonts,
        which have every Vietnamese letter. A test page in \code{acmart} had no missing letters
        and no warnings, with babel's Vietnamese captions (\vi{Hình~1}, \vi{Bảng~1},
        \vi{Phụ lục}), headings that \code{acmart} sets in capitals (\vi{GIỚI THIỆU}), and
        decimal commas in formulas (\vi{$-0,034$}, \vi{$p<0,001$}) through the package
        \code{icomma}. We set the other names ourselves: \vi{Tóm tắt} (babel's is \vi{Tóm tắt
        nội dung}), \vi{Từ khóa} and \vi{Lời cảm ơn}, which \code{acmart} has no Vietnamese
        for, and \vi{Tài liệu tham khảo} for the reference list, which babel calls only
        \vi{Tài liệu}.
  \item \textbf{The layout carries over.} The English \code{paper_v2}, built with LuaLaTeX and
        one added line (\code{fontspec}'s \code{quiet} option, which keeps font-feature notes
        out of the warnings), gives 14 pages and 0 warnings, with the code statement at the same
        place on page~8 and the references from the top of page~9, as with pdfLaTeX. So L1 starts
        from \code{paper_v2}'s exact layout, and any later change in length comes from the
        translation.
  \item \textbf{Overleaf.} The compiler and the main document are settings of the whole project
        (Menu). For \code{paper_v2_vn}, choose LuaLaTeX and \code{paper_v2_vn/paper_v2_vn.tex};
        for \code{paper/} and \code{paper_v2/}, go back to pdfLaTeX. Compiled with pdfLaTeX by
        mistake, \code{paper_v2_vn.tex} stops at once with a message saying so. If the switching
        is a nuisance, a separate Overleaf project with the three \code{.bib} files works the
        same.
  \item \textbf{Build time.} A LuaLaTeX pass took about 6~seconds here, against 4.7 for pdfLaTeX.
        If Overleaf's time limit ever becomes a problem, the four figures can be compiled once
        into PDF files and included as images.
  \item \textbf{Line breaks.} Vietnamese words are not hyphenated, which leaves the narrow
        columns less room. L1 allows a little more stretch, and each part fixes any line that
        still overflows.
\end{itemize}

\section{What L1 set up (10 November)}\label{sec:l1}

\begin{itemize}
  \item \textbf{The main file,} \code{paper_v2_vn.tex}, is \code{paper_v2.tex} with: the LuaLaTeX
        check (under pdfLaTeX the run stops after about a second with the one error
        ``paper\_v2\_vn.tex needs LuaLaTeX. In Overleaf: Menu, Compiler, LuaLaTeX.'' and no PDF);
        babel with Vietnamese as the main language and English for the rest; the names
        \vi{Tóm tắt}, \vi{Từ khóa}, \vi{Lời cảm ơn}, \vi{Tài liệu tham khảo} and \vi{Phân loại
        CCS} (Q7); decimal commas in formulas (\code{icomma}) and on the figures' axes
        (pgfplots' \code{use comma}); 1\,em of emergency stretch; Vietnamese as the PDF's
        language; and the names of N6 in the author block, the acknowledgments and the PDF's
        author field. The CCS heading needs a small patch of \code{acmart}'s \verb|\maketitle|,
        whose CCS block is the same in \code{acmart} 2.03 (2024) and 2.21 (2026); should it ever
        fail, the heading stays English with a warning, not an error.
  \item \textbf{Text not yet translated:} each section file is read inside \code{entext}, which
        sets it as English (hyphenation and spacing) under the Vietnamese captions. Each part
        removes the wrapper from the files it translates. The reference list and ACM's CCS terms
        are set as English for good (N7). The tables and figures, translated in L3 but read in
        by sections still in English, each sit in \code{vitext}, which sets them as Vietnamese;
        L9 removes these wrappers once no section is in \code{entext}. (L9 removed both kinds of
        wrapper.)
  \item \textbf{\code{numbers.tex}} is written by \code{make_numbers.py --lang vi}: 971 of the
        1,262 values change. The other 291 are whole numbers below 1,000, the path sets, the
        commits and branches, the split fingerprint, the GPU's name, the names of two sets
        (\code{r6d2}, \code{m5}) and the five version numbers. The figure data are identical in
        both languages. Without \code{--lang vi}, the script writes the same \code{numbers.tex}
        as before apart from the date in its second line. Until the author committed the script,
        the file recorded the LMSPS commit \code{0bf5bc1}; rebuilt from the commit, \code{5e32eb6},
        it records that, and no value changed. Run from \code{5e32eb6} without the option, the
        script writes \code{paper_v2}'s \code{numbers.tex} apart from the date and the commit, so
        \code{0bf5bc1} still stands for \code{paper/} and \code{paper_v2/}.
  \item \textbf{Checks:} 0 errors, 0 warnings, no overfull or underfull line and no missing
        letter; 14 pages. Pages 1--9 end at the same words as \code{paper_v2}'s, with the code
        statement at the same place on page~8 and the references from the top of page~9; in the
        appendix, pages 10 and 12 end a line or two earlier, a side effect of small changes in
        width such as \vi{Bảng} for Table in the captions. BibTeX's 9 notices are \code{paper_v2}'s (V3 of \code{plan_v2}). A
        script compares each file's \verb|\val|, \verb|\cite|, \verb|\ref|, \verb|\label|,
        \verb|\input| and \verb|\code| keys with \code{paper_v2}'s: all match.
\end{itemize}

\section{The final checks (L9, 10 November)}\label{sec:l9}

\begin{itemize}
  \item \textbf{Wrappers:} the 15 \code{vitext} wrappers and the two helper environments are gone;
        the PDF has the same words in the same order on every page before and after.
  \item \textbf{Keys:} every file has \code{paper_v2}'s \verb|\val|, \verb|\cite|, \verb|\ref|,
        \verb|\label|, \verb|\input| and \verb|\code| keys; the one allowed addition is
        \code{--lang vi} in Appendix~G (L8).
  \item \textbf{English:} none is left outside the quotations (N13), the terms given in
        parentheses at their first mention and the node types' English names in Section~3
        (checked by script).
  \item \textbf{Glossary:} no term that the glossary rules out (such as \vi{tập kiểm tra} or
        \vi{người đánh giá}); every \vi{đánh giá} is an evaluation, every \vi{kiểm tra} a check and
        every \vi{kiểm định} a statistical test.
  \item \textbf{Typography:} no straight quote and no space before punctuation; a tie before every
        cross-reference and between every number and its unit (two units and one awkward
        \vi{được được} fixed in L9); headings wrap only where a tie chooses the break (5.1 and
        Appendix~E).
  \item \textbf{Numbers:} \code{numbers.tex} and the 11 figure data files are byte for byte what
        \code{make_numbers.py --lang vi} writes.
  \item \textbf{PDF:} the title, author, keywords and language (\code{vi}) are in its XMP
        metadata; as in \code{paper_v2}, its older document-information block lacks the author and
        keywords, which \code{acmart}'s metadata package leaves out.
  \item \textbf{Build and length:} LuaLaTeX, 0 errors, 0 warnings, no overfull or underfull line
        and no missing letter; BibTeX's 9 notices as in \code{paper_v2}. 16 pages
        (\code{paper_v2}: 14). The code statement starts 8.85 pages in and the references 9.13
        (\code{paper_v2}: 7.77 and 8.02), so the main text is about 1.1 pages longer, within N15.
\end{itemize}

\section{Decisions}\label{sec:decisions}

\begin{enumerate}[label=N\arabic*]
  \item \textbf{Frozen versions (10 November):} \code{paper/} and \code{paper_v2/} change only
        for a must (Section~\ref{sec:rules}).
  \item \textbf{The folder (10 November):} \code{paper_v2_vn/}, which the author creates in the
        Overleaf project.
  \item \textbf{Purpose and length (10 November):} the same format as \code{paper_v2},
        translated in full. A main text up to about one page longer than \code{paper_v2}'s
        8 pages is fine; if it runs clearly longer, L9 reports it and asks before anything is
        cut.
  \item \textbf{Numbers (10 November):} the Vietnamese way, with a decimal comma and a dot
        between thousands (\vi{0,666}; \vi{1.147.920}; \vi{82,2\%}). Version numbers, commits,
        branches, the split fingerprint, the GPU's name and path names stay as they are.
        \code{make_numbers.py} gets a \code{--lang vi} option, which I write in L1 and the author
        commits. It writes the Vietnamese \code{numbers.tex}: these separators, and Vietnamese
        forms for the nine values that hold English words, which are seven lists of places such
        as ``6th, 7th and 8th'' (\vi{thứ 6, 7 và 8}) and the two counts of the McAuley data such
        as ``142.8 million'' (\vi{142,8 triệu}). Without the option, the script writes what it
        wrote before (only the date in the file's second line changes, as on every run), so
        \code{paper/} and \code{paper_v2/} are untouched. The Vietnamese \code{numbers.tex}
        records the new commit of the LMSPS repository.
  \item \textbf{Terms (10 November):} no Vietnamese terms exist yet for this work; L2 proposes
        the glossary for the author's approval.
  \item \textbf{Names (10 November):} the author, \vi{Phan Hoàng Tuấn Anh}; the supervisors in
        the acknowledgments, \vi{TS.~Ngô Dương Hà} and \vi{TS.~Phạm Nguyễn Huy Phương}; the
        affiliation, \vi{Khoa Công nghệ Thông tin, Trường Đại học Công Thương Thành phố Hồ Chí
        Minh (HUIT)}.
  \item \textbf{References (10 November):} the reference list stays in English, as the works were
        published, and citations keep ``et al.''\ (``McAuley et al.~[14]''), so ACM's
        bibliography style is used unchanged.
  \item \textbf{Appendix (10 November):} Appendices A--G are translated too (L8).
  \item \textbf{CCS concepts (10 November):} the heading is \vi{Phân loại CCS}, and ACM's terms
        stay in English (Q7).
  \item \textbf{Meta-path (10 November):} \vi{đường dẫn meta}, as proposed, with ``(meta-path)''
        at the first mention, and \vi{đường dẫn} for a path (Q8: ``leave as is'').
  \item \textbf{Validation and test (10 November):} \vi{xác thực} and \vi{thử nghiệm} (Q9). So
        \vi{kiểm định} stays free for statistical tests (\vi{kiểm định t} of Welch), an
        experiment is \vi{thực nghiệm} and a check \vi{kiểm tra}.
  \item \textbf{Seed and epoch (10 November):} the English words stay (Q10).
  \item \textbf{Quotations (10 November):} the cited papers' English words stay, in quotation
        marks, followed by a Vietnamese translation in parentheses and the citation (Q11).
  \item \textbf{Review (10 November):} a product review is \vi{nhận xét}, a reviewer \vi{người
        nhận xét}, co-review \vi{đồng nhận xét} and a rating \vi{số sao} (Q12). \vi{Đánh giá}
        then means evaluation only.
  \item \textbf{Length (10 November):} the main text may end about 1.3 pages after
        \code{paper_v2}'s, as projected after L4 (Q13), instead of the one page of N3: the
        remaining parts keep natural wording rather than a terser one, and nothing is cut. L9
        reports the final length.
  \item \textbf{Acknowledgments (10 November):} the plain wording, as in the English: \vi{Tác giả
        xin cảm ơn TS.~Ngô Dương Hà và TS.~Phạm Nguyễn Huy Phương đã hướng dẫn} (Q14).
  \item \textbf{Title and meta-path (L10, 20 November):} the title is now \vi{Đánh giá chi phí và
        hiệu quả của việc tìm kiếm meta-path cho gợi ý sản phẩm dựa trên HAN}, a descriptive title
        chosen by the author instead of a translation of the English question. The subtitle is
        unchanged. \vi{meta-path} replaces \vi{đường dẫn meta} in all 54 places, with no gloss at
        the first mention, and the short form \vi{đường dẫn} for a path in all 266 places, so the
        paper uses one term. This supersedes N10. \code{paper/} and \code{paper_v2/} are not affected.
  \item \textbf{Abstract (L10, 20 November):} the author's reviewed abstract. It leaves out the clause on
        merging duplicate editions (with \verb|\val{before-testleak-pct}|) and the sentence on GAT beating
        both meta-path models, which the English abstract keeps. The searched sets are \vi{tổng hợp từ
        30 lần chạy}, and the last sentence reads \vi{chưa cho thấy hiệu quả vượt trội}.
  \item \textbf{CCS concepts (L10, 20 November):} removed, since the paper is not for an ACM venue
        (N9 is void). The keywords stay. The heading code for \vi{Phân loại CCS} stays in the
        preamble, so putting back the two \verb|\ccsdesc| lines restores the block.
  \item \textbf{Punctuation (L10, 20 November):} no semicolons and no comma before \vi{và}, in the
        text, the captions, the tables, the figures and the PDF keywords (rule S10 of the glossary).
        All 155 semicolons and 94 such commas are gone. Most became a full stop or were dropped. A
        semicolon inside parentheses became a comma, and in a table cell a comma or a full stop,
        never a full stop before a number, since thousands take a dot (15.831). Lists whose items
        contain commas are numbered (1), (2). Where the plain change read badly or was ambiguous,
        the sentence was rewritten with the same content: the three ways and the three reasons in
        the introduction, the manual sets' four principles, the future work, the parts of the
        appendix that list files, and the AI statement. No \verb|\val|, citation or reference
        changed, and the main text ends 0.05 page later (9.11 pages in).
  \item \textbf{Introduction (L10, 20 November):} two changes by the author. The first of the three
        reasons is clearer: \vi{Thứ nhất, LMSPS đánh giá meta-path bằng mạng riêng mà nó dùng khi
        tìm kiếm chứ không phải bằng HAN, nên những meta-path tốt nhất với mạng đó chưa chắc là tốt
        nhất với HAN.} The sentence on cleaning the data at the end of the third paragraph is
        removed, since the first contribution gives the same figure (\verb|\val{before-testleak-pct}|)
        and the one after merging. Section~1 is 0.04 page shorter. \code{paper_v2} keeps its
        English wording.
  \item \textbf{Random sets and the author's introduction (L10, 20 November):} the author's term
        \vi{tập chọn ngẫu nhiên} replaces \vi{tập rút ngẫu nhiên} everywhere: the abstract, the text,
        the tables, the figures and an appendix heading. The short form is \vi{tập ngẫu nhiên}, one
        draw is \vi{lần chọn} (\vi{6 meta-path, lần chọn 2}) and \vi{rút đều} is \vi{chọn đều}. So
        that \vi{chọn ngẫu nhiên} keeps one meaning, other sampling is \vi{lấy mẫu} (LMSPS's
        meta-paths, the negatives), ties are \vi{bằng điểm thì xếp ngẫu nhiên}, and the
        validation-selected set is \vi{vốn là một tập ngẫu nhiên}. \vi{SVD rút gọn} (truncated SVD)
        is not a draw and stays. The author's final introduction (20 November) drops \vi{một lần
        (one-shot)} from LMSPS's description, the gloss \vi{(leakage)}, and the sentence on the
        cheapest option, which Section~8 keeps. Its two limits are reworded, and its \vi{tác vụ}
        became \vi{bài toán}, as everywhere else. The main text ends 0.03 page later (9.11 pages in).
  \item \textbf{Embedding and path weights (L10, 20 November):} the author's choices.
        \vi{embedding} replaces \vi{vectơ nhúng} (8 places, with its first-mention gloss) and
        \vi{vectơ} where it is a learned user or product vector (16 places): \vi{embedding tự do},
        \vi{các embedding người dùng}, \vi{Embedding $h_p$}. The LSA title vectors and HAN's
        attention vector stay \vi{vectơ}, since they are an input and a parameter, not a learned
        representation. \vi{cân trọng số} becomes \vi{gán trọng số cho} (HAN in Section~1, MCRec in
        Section~2), the verb the results use (\vi{trọng số mà HAN gán cho nó}), rather than
        \vi{đo trọng số}, since attention assigns the weights and does not measure them.
  \item \textbf{Path instances (L10, 20 November):} the author found \vi{thể hiện} unclear, so an
        instance of a meta-path is \vi{một đường đi theo meta-path đó} (Section~3). This also ties
        it to \vi{số đường đi}, the walk count by which HAN's graphs are pruned (Section~5).
  \item \textbf{Input features (L10, 20 November):} the author found \vi{mã one-hot} and \vi{mã
        băm} unclear. As inputs they are vectors (N23): \vi{vectơ one-hot} of the product group,
        \vi{vectơ multi-hot} of the category-name tokens, and for a user \vi{một vectơ nhị phân
        64 chiều sinh từ ID khách hàng bằng hàm băm} (8 of its 64 positions set, as
        \code{hgb/dat_generator.py} builds it), in Section~4 and Table~5.
  \item \textbf{No Discussion or Threats section (L10, 20 November):} the author's choice, so
        \vi{Kết luận} is now Section~8. Their main points moved. Why the searched sets lose, with
        the caveat that this explanation is untested, ends Section~7.4, and why GAT leads is in
        Section~7.5. The conclusion gained the cheapest choice, the lesson that a search should be
        judged by the recommender it serves, and a \vi{Hạn chế} paragraph (the small schema, the
        single random split, the untuned configuration, the smaller search, the AI-assisted manual
        sets). The sampled-metric item is covered in Section~6. The training details and the
        leaf-only categories are Appendix~H. The cross-references in Sections 1, 2 and 4 follow.
        The main text ends 0.34 page earlier (8.74 pages in). \code{paper_v2} keeps its Sections 8
        and 9.
  \item \textbf{Related work wording (L10, 20 November):} one-shot search is described rather than
        named (\vi{chỉ huấn luyện một lần: mọi ứng viên được huấn luyện cùng lúc trong một siêu
        mạng}), since \vi{tìm kiếm một lần} or \vi{chạy một lần} can be read as a search run only
        once. Yu et al.'s finding is
        paraphrased with its citation instead of quoted, and RMS's random baseline is named
        \vi{đối chứng ngẫu nhiên}. A top-k query finds \vi{$k$ meta-path tốt nhất}, as Zhu et al.
        write.
  \item \textbf{AI statement (L10, 20 November):} the author's final wording. It does not mention
        the drafting, editing or translation of the text, and it says the AI \vi{hỗ trợ viết} the
        script. The assistant told the author that this under-reports its part (it drafted and
        translated the text and wrote the script) and proposed wording that names the English
        version or not. The author kept this wording.
  \item \textbf{Search against manual design (L10, 21 November):} the author's framing. The paper
        compares the searched sets with the manual sets, and the random sets are the control
        (\vi{đối chứng}) that shows how much the choice of meta-paths matters. Changed: the abstract
        (two ways plus a control, the last sentence on search against manual design), the keywords
        (\vi{thiết kế meta-path thủ công} for \vi{đối chứng ngẫu nhiên}), the introduction (the
        question, two ways plus a control, the findings, contributions 3 and 4), the related-work
        paragraph (now \vi{Đánh giá phương pháp tìm kiếm}, starting from LMSPS's HAN comparison with
        HAN's original meta-paths), the setup list, Sections~7.2 (\vi{Đối chứng ngẫu nhiên}, now with
        the cheapest choice) and 7.3 (\vi{Chi phí so với hiệu quả}), the conclusion and the overview
        figure. GAT and the popularity model are \vi{mô hình cơ sở}. No number, table or result
        changed, and the random sets' result is still in the abstract. \code{paper_v2} keeps its
        framing.
  \item \textbf{A verdict that depends on the context (L10, 22 November):} the author's framing. The
        searched sets come fairly close to the manual ones (\vi{khá gần}, not \vi{rất gần}, as the
        gap on the pre-registered metric is significant), so the choice depends on the context.
        Manual design is better when an expert can design the meta-paths in less time than the
        search takes. Without one, search gives nearly equivalent sets (\vi{gần tương đương}) with no
        domain knowledge. It is not written that search wins then, since the random sets did at
        least as well at no cost. The abstract, the introduction and the conclusion say this, and
        Section~7.3 gives the yardstick (about 20 GPU hours). The conclusion's limitations add that
        the manual sets were not designed by a domain expert.
  \item \textbf{No appendix, a supplementary document (L10, 22 November):} the author's choice. The
        paper ends with the references, at ten pages. The appendix became a document of its own,
        \code{supplement_vn.tex} (\vi{Tài liệu bổ sung}, seven pages), which the author keeps ready
        in case someone asks for the details, so the paper does not refer to it. Moved into the
        paper: the ten random sets (rows of the sets table, the 4- and 6-path draws side by side so
        that the paper stays at ten pages), the key settings and the software versions
        (Section~6), why Recall@10 and Precision@10 are not reported, the numbers behind four claims
        of Section~7 (the user groups, HAN against metapath2vec, the random sets' weak link to the
        search's preference, GAT's lead by user group) and the runs' agreement (Section~5), the two
        minor limitations (one sentence each) and, in the code statement, the commits, the data
        rebuild and the records. The quotation in Section~4 now has its source, the McAuley data
        page, as a footnote, which was in the appendix. Only in the supplementary document: the
        overview figure, the tables of the graph, the changes, the settings, the user groups and
        the full results, the search's details and the commands. Its sections are lettered A--H
        and called \vi{Mục}, and its tables and figure are numbered S1, S2, \ldots. It prints the
        paper's numbers through the \code{xr} package (labels prefixed \code{P-}, printed with
        \verb|\ref*| so that no link points into the other PDF), so the paper is built first:
        \code{latexmkrc} does this with Overleaf's rule for \code{xr}. Column balancing stays off,
        since acmart's \verb|\balance| would come in the second column of the last page and warn.
  \item \textbf{Ancestor categories are \vi{danh mục cấp trên} (L10, 22 November):} the author's
        choice, for \vi{tổ tiên}. The author first wrote \vi{danh mục cha}, which is only the direct
        parent, while the numbers (48,7\% of the products in one category, a median of 13
        categories) are for every category up to the root. Changed: Section~4, the conclusion's
        limitations, two places in the supplementary document and the glossary.
\end{enumerate}

\section{Questions}\label{sec:questions}

Q1--Q7 were answered on 10 November (N3--N9), and Q8--Q12, the five term choices of L2, the same
day (N10--N14), and Q13, on the length, after L4 (N15). The rest of the glossary and the sample
stand as proposed unless the author marks changes. Q14, on the acknowledgments, was answered
after L7 (N16). None is open.

\section{The parts}\label{sec:parts}

{\small
\begin{tabularx}{\textwidth}{@{}>{\raggedright\arraybackslash}p{0.2\textwidth}L>{\raggedright\arraybackslash}p{0.13\textwidth}@{}}
  \toprule
  Part & What it does & Needs \\
  \midrule
  L1 Set up & \code{paper_v2_vn/} as a copy of \code{paper_v2} (still in English) that builds with
    LuaLaTeX: the Vietnamese language setup (captions, the names of the abstract, keywords,
    acknowledgments and CCS concepts, decimal commas in formulas and figure axes, the pdfLaTeX
    stop, a little more stretch), the names with diacritics (N6), and \code{numbers.tex} in the
    format of N4, with \code{make_numbers.py} for you to commit. Checks: 0 warnings and
    \code{paper_v2}'s layout. Sends the whole new folder and the script. Done 10 November
    (Section~\ref{sec:l1}). & N4, N6 \\
  L2 Glossary and style & A glossary of the paper's terms (English, Vietnamese, the form used at
    first mention), proposed for your approval (N5); the style rules (\vi{chúng tôi} for ``we'', sentence
    case in headings, \vi{Mục}, \vi{Bảng}, \vi{Hình} and \vi{Phụ lục} for Section, Table, Figure
    and Appendix, spaces and punctuation); and, as a sample of the style, the title, subtitle,
    abstract and keywords. For your approval before L3. Done 10 November: 131 entries and 12
    rules in \code{glossary_vn.tex}, updated for N10--N14; the Vietnamese abstract moves the main
    text's end about 0.06 page later. & N5, N7 \\
  L3 Tables and figures & The 11 tables and 4 figures: captions, column heads, notes, labels,
    legends, the pipeline diagram and the axes; every width checked so that nothing overflows.
    Done 10 November: also the figures' descriptions for screen readers; the base papers' own
    settings in Table~6 in the Vietnamese number format (0,005; 1.000); the main text ends about
    0.12 page later (7.95 pages in). & L2 \\
  L4 Sections 1--3 & Introduction, Related work, Preliminaries (about 1,300 English words). Done
    10 November: 0.24 page longer than the English; the node types' English names added in
    Section~3 to explain the letters P, C and U; eight new terms in the glossary. & L2 \\
  L5 Sections 4--5 & Data and Method (about 1,900 words). Done 10 November: 0.33 page longer
    than the English, so the main text ends 8.52 pages in, on course for the 9.0 of N15; five
    new terms in the glossary. & L2 \\
  L6 Sections 6--7 & Setup and Results (about 1,700 words and 118 numbers). Done 10 November:
    0.19 page longer than the English, so the main text ends 8.71 pages in; at the rate so far,
    Sections 8--10 and the end statements bring it to about 8.9; the document has 16 pages. & L2
    \\
  L7 Sections 8--10 and the end & Discussion, Threats to validity, Conclusion, the code and data
    statement, the AI statement with the translation added, and the acknowledgments (about 1,050
    words). Done 10 November: the AI statement adds that the assistant translated this version,
    following a glossary that it proposed and the author decided, and that the author reviewed
    the translation; the code statement starts 8.85 pages in (\code{paper_v2}: 7.77), so the
    main text is about 1.1 pages longer, within N15, and the references begin 9.13 pages in. & L2
    \\
  L8 Appendix & Appendices A--G (about 1,700 words of prose; their tables come from L3). Done
    10 November: one deliberate addition to \code{paper_v2}, in Appendix~G's Records, which
    names \code{make_numbers.py}'s \code{--lang vi} option because it writes this version's
    numbers (the parity script lists it as allowed, with that reason); the heading of
    Appendix~E tied so that it breaks before \vi{và}; 16 pages in all. & L2, N8 \\
  L9 Final pass & The whole paper against \code{paper_v2}: every key, table row and figure label;
    the \code{vitext} wrappers removed;
    the glossary applied throughout; no English left by mistake; typography; the PDF's title,
    keywords and language; the page count against N3. Sends the complete folder. Done 10
    November (Section~\ref{sec:l9}). & L3--L8, N3 \\
  \bottomrule
\end{tabularx}\par}

If the author or the supervisors send corrections after reading the Vietnamese text, they come
in as one more part, L10, after L9 or after any earlier part.

\end{document}
